% Part: first-order-logic
% Chapter: syntax-and-semantics
% Section: formation-sequences

\documentclass[../../../include/open-logic-section]{subfiles}

\begin{document}

\olfileid{fol}{syn}{fseq}

\olsection{تشکیل دے سلسلے}

!!{formula}s نوں استقرائی تعریف نال تعریف کرنا، تے !!{formula}s
دیاں خاصیتاں استقرا نال ثابت کرن دا نال دا طریقہ، اک سوہنا تے
کارگر ڈھنگ اے۔ پر !!{formula}s دی بناوٹ نوں ہِٹھوں اُتّے ول،
قدم بہ قدم وی ویکھیا جا سکدا اے۔ ایتھے اسیں ایہ کم
\emph{تشکیل دے سلسلے} دے تصور نال کردے آں۔
%
ایس طریقے نال ٹرم تے !!{formula}s نوں اوہناں دیاں استقرائی
تعریفاں دا حوالہ دتے بغیر پیش کرن لئی، پہلے اسیں کسے
زبان~$\Lang L$ دیاں علامتاں توں بنی من مانی لڑی دا تصور دیندے آں۔

\begin{defn}[لڑیاں]
\ollabel{defn:string}
فرض کرو $\Lang L$ پہلے درجے دی زبان اے۔ \emph{$\Lang
L$-لڑی}، $\Lang L$ دیاں علامتاں دا محدود سلسلہ اے۔ جدوں
سیاق توں زبان~$\Lang L$ صاف ہووے، تاں اسیں
$\Lang L$-لڑی نوں صرف \emph{لڑی} وی آکھاں گے۔
\end{defn}

\begin{ex}
کسے وی پہلے درجے دی زبان $\Lang L$ لئی، سارے
$\Lang L$-!!{formula}s، $\Lang L$-لڑیاں نیں؛ پر
لازم نہیں کہ ہر ایسی لڑی فارمولا ہووے۔ مثال لئی،
\[)(\Obj v_0\lif\lexists\] اک $\Lang L$-لڑی اے،
پر $\Lang L$-!!{formula} نہیں۔
\end{ex}

\begin{defn}[ٹرماں لئی تشکیل دے سلسلے]
\ollabel{defn:fseq-trm}
$\Lang L$-لڑیاں دا محدود سلسلہ $\tuple{t_0,\dotsc,t_n}$،
ٹرم $t$ لئی \emph{تشکیل دا سلسلہ} اے جے $t \ident t_n$،
تے ہر $i \leq n$ لئی یا تاں $t_i$ !!a{variable} یا
!!a{constant} اے، یا $\Lang L$ وچ اک $k$-مقامی
!!{function}~$f$ اے تے ایہو جہے $m_1,\dotsc,m_k < i$
موجود نیں کہ $t_i \ident f(t_{m_1},\dotsc,t_{m_k})$۔
% PNB-SYN-003: منجمد ماخذ وچ k-مقامی فنکشن لئی m_0 توں m_k تک k+1 دلیلاں سن؛ ایتھے k دلیلاں رکھیاں نیں۔
جدوں فرق واضح کرن دی لوڑ ہووے، اسیں ٹرماں دے تشکیل دے سلسلیاں نوں
\emph{ٹرم تشکیل دے سلسلے} آکھاں گے۔
\end{defn}

\begin{ex}
ایہ سلسلہ
\[
    \tuple{\Obj c_0, \Obj v_0, \Atom{\Obj f^2_0}{\Obj c_0, \Obj v_0}, \Atom{\Obj f^1_0}{\Atom{\Obj f^2_0}{\Obj c_0, \Obj v_0}}}
\]
ٹرم $\Atom{\Obj f^1_0}{\Atom{\Obj
f^2_0}{\Obj c_0, \Obj v_0}}$ لئی تشکیل دا سلسلہ اے،
تے ایہ وی:
\[
    \tuple{\Obj v_0, \Obj c_0, \Atom{\Obj f^2_0}{\Obj c_0, \Obj v_0}, \Atom{\Obj f^1_0}{\Atom{\Obj f^2_0}{\Obj c_0, \Obj v_0}}}.
\]
\end{ex}

\begin{defn}[فارمولیاں لئی تشکیل دے سلسلے]
\ollabel{defn:fseq-frm}
$\Lang L$-لڑیاں دا محدود سلسلہ $\tuple{!A_0,\dotsc,!A_n}$
$!A$ لئی \emph{تشکیل دا سلسلہ} اے جے $!A \ident !A_n$،
تے ہر $i \leq n$ لئی یا تاں $!A_i$ ایٹمی !!{formula} اے،
یا ایہو جہے $j,k < i$ تے !!a{variable}~$x$ موجود نیں
کہ اگلیاں وچوں کوئی اک صورت پوری ہُندی اے:
\begin{enumerate}
    \tagitem{prvNot}{$!A_i \ident \lnot !A_j$.}{}%
    \tagitem{prvAnd}{$!A_i \ident (!A_j \land !A_k)$.}{}%
    \tagitem{prvOr}{$!A_i \ident (!A_j \lor !A_k)$.}{}%
    \tagitem{prvIf}{$!A_i \ident (!A_j \lif !A_k)$.}{}%
    \tagitem{prvIff}{$!A_i \ident (!A_j \liff !A_k)$.}{}%
    \tagitem{prvAll}{$!A_i \ident \lforall[x][!A_j]$.}{}%
    \tagitem{prvEx}{$!A_i \ident \lexists[x][!A_j]$.}{}%
\end{enumerate}
جدوں فرق واضح کرن دی لوڑ ہووے، اسیں فارمولیاں دے تشکیل دے سلسلیاں نوں
\emph{فارمولا تشکیل دے سلسلے} آکھاں گے۔
\end{defn}

\begin{ex}
\[
\tuple{
    \Atom{\Obj A^1_0}{\Obj v_0},
    \Atom{\Obj A^1_1}{\Obj c_1},
    (\Atom{\Obj A^1_1}{\Obj c_1} \land \Atom{\Obj A^1_0}{\Obj v_0}),
    \lexists[\Obj v_0][(\Atom{\Obj A^1_1}{\Obj c_1} \land \Atom{\Obj A^1_0}{\Obj v_0})]
}
\]
$\lexists[\Obj v_0][(\Atom{\Obj A^1_1}{\Obj
c_1} \land \Atom{\Obj A^1_0}{\Obj v_0})]$ لئی
تشکیل دا سلسلہ اے؛ ایہ وی اک اے:
\begin{multline*}
\tuple{
    \Atom{\Obj A^1_0}{\Obj v_0},
    \Atom{\Obj A^1_1}{\Obj c_1},
    (\Atom{\Obj A^1_1}{\Obj c_1} \land \Atom{\Obj A^1_0}{\Obj v_0}),
    \Atom{\Obj A^1_1}{\Obj c_1},\\
    \lforall[\Obj v_1][\Atom{\Obj A^1_0}{\Obj v_0}],
    \lexists[\Obj v_0][(\Atom{\Obj A^1_1}{\Obj c_1} \land \Atom{\Obj A^1_0}{\Obj v_0})]
}.
\end{multline*}
%
دوجی مثال وکھاؤندی اے کہ تشکیل دے سلسلیاں وچ ”فالتو مواد“ وی
آ سکدا اے: اوہ !!{formula}s جیہڑے دہرائے گئے ہون یا
بناوٹ وچ کوئی حصہ نہ پاؤندے ہون۔
\end{ex}

\begin{prop}\ollabel{prop:formed}
ہر !!{formula}~$!A$ جیہڑا~$\Frm[L]$ وچ اے،
اوہدے لئی تشکیل دا سلسلہ موجود اے۔
\end{prop}

\begin{proof}
فرض کرو $!A$ ایٹمی اے۔ تاں سلسلہ $\tuple{!A}$،
$!A$ لئی تشکیل دا سلسلہ اے۔
%
ہُن فرض کرو $!B$ تے~$!C$ دے تشکیل دے سلسلے
ترتیب وار $\tuple{!B_0,\dotsc,!B_n}$ تے
$\tuple{!C_0,\dotsc,!C_m}$ نیں۔
%
\begin{enumerate}
  \tagitem{prvNot}{جے $!A \ident \lnot !B$،
      تاں $\tuple{!B_0,\dotsc,!B_n,\lnot !B_n}$
      $!A$ لئی تشکیل دا سلسلہ اے۔}{}
  \tagitem{prvAnd}{جے $!A \ident (!B \land !C)$،
      تاں $\tuple{!B_0,\dotsc,!B_n,!C_0,\dotsc,!C_m,(!B_n \land !C_m)}$
      $!A$ لئی تشکیل دا سلسلہ اے۔}{}
  \tagitem{prvOr}{جے $!A \ident (!B \lor !C)$،
      تاں $\tuple{!B_0,\dotsc,!B_n,!C_0,\dotsc,!C_m,(!B_n \lor !C_m)}$
      $!A$ لئی تشکیل دا سلسلہ اے۔}{}
  \tagitem{prvIf}{جے $!A \ident (!B \lif !C)$،
      تاں $\tuple{!B_0,\dotsc,!B_n,!C_0,\dotsc,!C_m,(!B_n \lif !C_m)}$
      $!A$ لئی تشکیل دا سلسلہ اے۔}{}
  \tagitem{prvIff}{جے $!A \ident (!B \liff !C)$،
      تاں $\tuple{!B_0,\dotsc,!B_n,!C_0,\dotsc,!C_m,(!B_n \liff !C_m)}$
      $!A$ لئی تشکیل دا سلسلہ اے۔}{}
  \tagitem{prvAll}{جے $!A \ident \lforall[x][!B]$،
      تاں $\tuple{!B_0,\dotsc,!B_n,\lforall[x][!B_n]}$
      $!A$ لئی تشکیل دا سلسلہ اے۔}{}
  \tagitem{prvEx}{جے $!A \ident \lexists[x][!B]$،
      تاں $\tuple{!B_0,\dotsc,!B_n,\lexists[x][!B_n]}$
      $!A$ لئی تشکیل دا سلسلہ اے۔}{}
\end{enumerate}
!!{formula}s اُتے استقرا دے اصول موجب ہر !!{formula}
دا تشکیل دا سلسلہ اے۔
\end{proof}

اسیں اُلٹا رخ وی ثابت کر سکدے آں۔ ایہ اہم اے،
کیوں جے ایہ وکھاؤندا اے کہ فارمولیاں نوں تعریف
کرن دے ساڈے دونیں طریقے ہم ارز نیں: اوہ اکّو
نتیجہ دیندے نیں۔ ایس دا مطلب ایہ وی اے کہ
اسیں تشکیل دے سلسلے دی لمبائی اُتے عام استقرا
لا کے فارمولیاں بارے قضیے ثابت کر سکدے آں۔

\begin{lem}
\ollabel{lem:fseq-init}
فرض کرو $\tuple{!A_0,\dotsc,!A_n}$،~$!A_n$ لئی
تشکیل دا سلسلہ اے، تے $k \leq n$۔ تاں
$\tuple{!A_0,\dotsc,!A_k}$،~$!A_k$ لئی
تشکیل دا سلسلہ اے۔
\end{lem}

\begin{proof}
مشق۔
\end{proof}

\begin{prob}
\olref[fol][syn][fseq]{lem:fseq-init} ثابت کرو۔
\end{prob}

\begin{thm}
\ollabel{thm:fseq-frm-equiv}
$\Frm[L]$ اوہناں ساریاں $\Lang L$-لڑیاں~$!A$
دا سیٹ اے جیہناں دی شرط ایہ اے کہ $!A$ لئی فارمولا تشکیل دا سلسلہ موجود اے۔
\end{thm}

\begin{proof}
$F$ نوں زبان~$\Lang L$ دیاں اوہناں ساریاں علامتی
لڑیاں دا سیٹ منّو جنہاں دا تشکیل دا سلسلہ اے۔
اسیں \olref[fol][syn][fseq]{prop:formed} وچ
ویکھ چکے آں کہ $\Frm[L] \subseteq F$؛ سو ہُن
اسیں اُلٹا رخ ثابت کردے آں۔

فرض کرو $!A$ دا تشکیل دا سلسلہ
$\tuple{!A_0,\dotsc,!A_n}$ اے۔ اسیں ثابت
کردے آں کہ $!A \in \Frm[L]$؛ ایہ کم
$n$ یعنی آخری رکن دے اشاریے اُتے مضبوط استقرا نال کردے آں۔
ساڈا استقرائی مفروضہ اے کہ جیہڑی وی علامتی لڑی
کسے ایسے تشکیل دے سلسلے دا آخری رکن اے
جیہدا آخری اشاریہ $m < n$ اے، اوہ $\Frm[L]$ وچ اے۔
% PNB-SYN-004: ماخذ n اُتے استقرا کردا پر n+1 رکناں نوں لمبائی n آکھدا اے؛ ایتھے آخری اشاریہ مستقل رکھیا اے۔
تشکیل دے سلسلے دی تعریف موجب یا تاں
$!A \ident !A_n$ ایٹمی اے، یا ایہو جہے
$j,k < n$ موجود نیں کہ اگلیاں وچوں کوئی اک صورت اے:
\begin{enumerate}
    \tagitem{prvNot}{$!A \ident \lnot !A_j$.}{}%
    \tagitem{prvAnd}{$!A \ident (!A_j \land !A_k)$.}{}%
    \tagitem{prvOr}{$!A \ident (!A_j \lor !A_k)$.}{}%
    \tagitem{prvIf}{$!A \ident (!A_j \lif !A_k)$.}{}%
    \tagitem{prvIff}{$!A \ident (!A_j \liff !A_k)$.}{}%
    \tagitem{prvAll}{$!A \ident \lforall[x][!A_j]$.}{}%
    \tagitem{prvEx}{$!A \ident \lexists[x][!A_j]$.}{}%
\end{enumerate}
ہُن اسیں صورتاں دے لحاظ نال دلیل دیندے آں۔
جے $!A$ ایٹمی اے تاں $!A_n \in \Frm[L]$۔
ایس دی بجائے فرض کرو کہ $!A \ident (!A_j \land !A_k)$۔
% PNB-SYN-005: ماخذ دے دو Frm[L_0] حوالیاں نوں پہلے درجے دے Frm[L] نال تے equiv نوں نحوی ident نال درست کیتا اے۔
\olref[fol][syn][fseq]{lem:fseq-init} موجب
$\tuple{!A_0,\dotsc,!A_j}$ تے
$\tuple{!A_0,\dotsc,!A_k}$ ترتیب وار
$!A_j$ تے~$!A_k$ لئی تشکیل دے سلسلے نیں۔
ایہ $!A$ دے سلسلے دے پورے توں چھوٹے مُڈھلے
ذیلی سلسلے نیں؛ اوہناں دے آخری اشاریے
اصل سلسلے دے آخری اشاریے~$n$ توں گھٹ نیں۔
ایس لئی استقرائی مفروضے
موجب $!A_j$ تے~$!A_k$،~$\Frm[L]$ وچ نیں،
تے !!a{formula} دی تعریف موجب
$(!A_j \land !A_k)$ وی اوہدے وچ اے۔
باقی صورتاں لئی وی ایسّے طرح دی دلیل اے۔
\end{proof}

ٹرماں لئی تشکیل دے سلسلیاں دیاں خاصیتاں
!!{formula}s دے سلسلیاں دیاں خاصیتاں ورگیاں نیں۔

\begin{prop}
\ollabel{prop:fseq-trm-equiv}
$\Trm[L]$ اوہناں ساریاں $\Lang L$-لڑیاں $t$
دا سیٹ اے جیہناں دی شرط ایہ اے کہ $t$ لئی ٹرم تشکیل دا سلسلہ موجود اے۔
\end{prop}

\begin{proof}
مشق۔
\end{proof}

\begin{prob}
\olref[fol][syn][fseq]{prop:fseq-trm-equiv} ثابت کرو۔
اشارہ: \olref[fol][syn][fseq]{thm:fseq-frm-equiv}
دے ثبوت ورگی حکمتِ عملی ورتو۔
\end{prob}

تشکیل دے سلسلیاں وچ دو طرح دا ”فالتو مواد“ آ
سکدا اے: دہرائے گئے رکن، تے اوہ رکن جیہڑے
فارمولے یا ٹرم دی بناوٹ لئی غیر متعلق ہون۔
اسیں گھٹ توں گھٹ تشکیل دے سلسلیاں نوں ویکھ
کے دونیں طرح دے فالتو رکناں نوں ہٹا سکدے آں۔

\begin{defn}[گھٹ توں گھٹ تشکیل دے سلسلے]
\ollabel{defn:minimal-fseq}
تشکیل دا سلسلہ $\tuple{!A_0, \ldots, !A_n}$،
جیہڑا فارمولا~$!A$ لئی اے،~$!A$ لئی
\emph{گھٹ توں گھٹ تشکیل دا سلسلہ} اے جے
ہر ہور تشکیل دے سلسلے~$s$، جیہڑا~$!A$ لئی ہووے،
نوں لے کے ایس~$s$ دی لمبائی~$n+1$ توں وڈّی
یا برابر ہووے۔

ایسّے طرح تشکیل دا سلسلہ $\tuple{t_0, \ldots, t_n}$،
جیہڑا ٹرم~$t$ لئی اے،~$t$ لئی
\emph{گھٹ توں گھٹ تشکیل دا سلسلہ} اے جے
ہر ہور تشکیل دے سلسلے~$s$، جیہڑا~$t$ لئی ہووے،
نوں لے کے ایس~$s$ دی لمبائی~$n+1$ توں وڈّی
یا برابر ہووے۔
\end{defn}

دھیان دیو کہ کسے فارمولے یا ٹرم لئی اک توں
ودھ گھٹ توں گھٹ تشکیل دے سلسلے ہو سکدے نیں،
پر اوہناں وچ عین اوہی لڑیاں شامل ہون گیاں۔

\begin{prop}
\ollabel{prop:subformula-equivs}
اگلیاں گلاں ہم ارز نیں:
\begin{enumerate}
    \item $!B$،~$!A$ دا ذیلی !!{formula} اے۔
    \item $!B$،~$!A$ دے ہر تشکیل دے سلسلے وچ آؤندا اے۔
    \item $!B$،~$!A$ دے کسے گھٹ توں گھٹ تشکیل دے سلسلے وچ آؤندا اے۔
\end{enumerate}
\end{prop}

\begin{proof}
مشق۔
\end{proof}

\begin{prob}
\olref[fol][syn][fseq]{prop:subformula-equivs} ثابت کرو۔
\end{prob}

\begin{history}
تشکیل دے سلسلے Raymond Smullyan نے اپنی
کتاب \emph{First-Order Logic} \citep{Smullyan1968}
وچ پیش کیتے سن۔ تشکیل دے سلسلیاں دیاں ہور خاصیتاں
\citet{Zuckerman1973} نے ثابت کیتیاں۔
\end{history}

\end{document}
